\section{Quadratic-field coefficients and a uniform signed average}\label{sec:quadratic}

A quadratic field has a useful feature: its zeta function contains ordinary $\zeta$ as a factor. Varying the field therefore gives many different arithmetic expressions with the same target factor. The question is whether averaging these expressions can simplify the auxiliary arithmetic without losing ordinary $\zeta$.

The first result replaces ideal coefficients by an explicit sequence on rational integers. The replacement preserves any cancellation estimate with exponent greater than $1/2$. The second result shows quantitatively that averaging this sequence over quadratic fields recovers the ordinary M\"obius sequence. Neither step proves that the resulting sum is small; it tells us exactly which sum must be estimated.

The fixed-field construction is from \cite{R996}, and the uniform extension and field average are from \cite{R1002}. The family here differs from the dyadic family in Section~\ref{sec:residuals}; we specify each one rather than interchange them.

\subsection{Replacing ideal coefficients by rational coefficients}
Put $\mu_6(n)=\mu(n)\ind_{(n,6)=1}$. For a negative fundamental discriminant $\Delta$, let $\chi_\Delta$ be the associated primitive quadratic character. Define
\begin{equation}\label{eq:primitive-sequences}
 a_\Delta=\mu_6*(\mu_6\chi_\Delta),\qquad
 b_\Delta(n)=\mu_6(n)\prod_{p\mid n}(1+\chi_\Delta(p)),
\end{equation}
where $*$ denotes Dirichlet convolution. Thus $b_\Delta$ is supported on squarefree integers coprime to $6$. In particular a ramified prime $p\ge5$ has coefficient $-1$, not zero. Write $\zeta_6$ and $L_6$ for deletion of the Euler factors at $2,3$.

\begin{theorem}[Uniform primitive-core comparison]\label{thm:primitive}
On $\Re s>1$, let
\[
 A_\Delta(s)=\sum_n a_\Delta(n)n^{-s}
 =\frac1{\zeta_6(s)L_6(s,\chi_\Delta)},\qquad
 B_\Delta(s)=\sum_n b_\Delta(n)n^{-s}.
\]
There are Dirichlet series $V_\Delta,W_\Delta$ such that
\begin{equation}\label{eq:primitive-inverse}
 B_\Delta=A_\Delta V_\Delta,\qquad
 A_\Delta=B_\Delta W_\Delta,\qquad V_\Delta W_\Delta=1.
\end{equation}
For every fixed $\sigma>1/2$, their coefficients satisfy
\[
 \sum_n\frac{\abs{v_\Delta(n)}}{n^\sigma}\le C_V(\sigma),\qquad
 \sum_n\frac{\abs{w_\Delta(n)}}{n^\sigma}\le C_W(\sigma),
\]
with constants independent of $\Delta$. Both functions are holomorphic and nonvanishing on $\Re s>1/2$.
\end{theorem}
\begin{proof}
At $p\ge5$ put $x=p^{-s}$. At a split prime the local factors are
\begin{align*}
 A_p(x)&=(1-x)^2,& B_p(x)&=1-2x,\\
 V_p(x)&=\frac{1-2x}{(1-x)^2}
       =1-\sum_{j\ge2}(j-1)x^j,&
 W_p(x)&=\frac{(1-x)^2}{1-2x}
       =1+\sum_{j\ge2}2^{j-2}x^j.
\end{align*}
At an inert prime they are
\[
 A_p=1-x^2,\quad B_p=1,\quad
 V_p=(1-x^2)^{-1},\quad W_p=1-x^2.
\]
At a ramified prime $A_p=B_p=1-x$, so $V_p=W_p=1$. All factors at $2,3$ are one. These are identities of complete power series; the higher prime powers are not discarded.

Since $2p^{-\sigma}<1$ for $p\ge5$ and $\sigma>1/2$, the products
\begin{align*}
 C_V(\sigma)&=\prod_{p\ge5}
  \left(1+\frac{p^{-2\sigma}}{(1-p^{-\sigma})^2}\right),\\
 C_W(\sigma)&=\prod_{p\ge5}
  \left(1+\frac{p^{-2\sigma}}{1-2p^{-\sigma}}\right)
\end{align*}
converge. They dominate the corresponding local absolute coefficient sums in every splitting case. This proves the uniform bounds and normal convergence on smaller half-planes. The local inverse identities pass to the products, proving nonvanishing and \eqref{eq:primitive-inverse}.
\end{proof}

\begin{corollary}[No loss in a fixed power exponent]\label{cor:primitive-power}
Fix $\beta>1/2$. A bound
\[
 \sum_{n\le X}a_\Delta(n)\ll_{\Delta,\eps}X^{\beta+\eps}
\]
for every $\eps>0$ is equivalent to the same bound with $b_\Delta$ in place of $a_\Delta$. The same equivalence holds for estimates uniform over additive phases, or over all real polynomial phases of any one fixed degree.
\end{corollary}
\begin{proof}
For example, $b_\Delta=a_\Delta*v_\Delta$ gives
\[
 \sum_{n\le X}b_\Delta(n)
 =\sum_{d\le X}v_\Delta(d)\sum_{m\le X/d}a_\Delta(m).
\]
Apply the hypothesized bound with a smaller positive loss and sum $\abs{v_\Delta(d)}d^{-\beta-\eps'}$. The inverse uses $W_\Delta$. For a polynomial phase $e^{2\pi\ii P(n)}$, the inner phase becomes $P(dm)$, which is still a polynomial of the same degree. The equivalence transfers an estimate; it does not establish one.
\end{proof}

At the all-$\eps$ square-root exponent, this reproduces the fixed-field reciprocal criterion described in \cite{R996}. For $K=\Q(\sqrt{-3})$, the unramified primitive sequence is simply $\mu(n)2^{\omega(n)}$ on squarefree integers all of whose prime factors are $1\bmod3$. Its fixed-field square-root estimate is stronger than ordinary RH alone, because it also controls the factor $L(s,\chi_{-3})$.

\subsection{Averaging over quadratic fields}
For $Q\ge2^{16}$ define
\[
 \mathcal D(Q)=\{d\le Q:d>0,\ d\text{ squarefree},\ d\equiv3\pmod4\},
 \qquad N_Q=\abs{\mathcal D(Q)}.
\]
Every $-d$ is a fundamental discriminant. We average uniformly over this complete set.

\begin{lemma}\label{lem:quadratic-mean}
One has $N_Q\ge Q/8$. For every odd squarefree $r>1$,
\begin{align}
 \left|\frac1{N_Q}\sum_{d\in\mathcal D(Q)}\chi_{-d}(r)\right|
 &\le24\sqrt{r/Q},\label{eq:mean-trivial}\\
 \left|\frac1{N_Q}\sum_{d\in\mathcal D(Q)}\chi_{-d}(r)\right|
 &\le\frac{24r^{1/4}\sqrt{1+\log r}}{\sqrt Q}.
 \label{eq:mean-pv}
\end{align}
\end{lemma}
\begin{proof}
There are at least $Q/4-1$ integers $3\bmod4$ up to $Q$. For each odd prime $p\le\sqrt Q$, at most $Q/(4p^2)+1$ are divisible by $p^2$. Using $\sum_{p\ge3}p^{-2}\le\sum_{n\ge3}n^{-2}\le4/9$ gives
\[
 N_Q\ge\frac{5Q}{36}-\sqrt Q-1\ge\frac Q8.
\]

Expand $\mu^2(d)=\sum_{a^2\mid d}\mu(a)$. Terms with $(a,r)>1$ vanish and are retained as zero terms. For the remaining odd $a$, the inner character sum is over a progression $4j+3$ of length at most $Q/a^2$. The Jacobi character modulo the odd squarefree integer $r$ is primitive and nonprincipal. Its complete mean is zero, so the inner sum is at most $\min(r,Q/a^2)$ in modulus. More generally, any bound $T$ for its incomplete sums gives
\[
 \sum_{a\le\sqrt Q}\min(T,Q/a^2)\le3\sqrt{QT}.
\]
This follows by splitting at $\sqrt{Q/T}$; when that point is below one, the bound follows directly from $\sum a^{-2}<2$ and $T>Q$.

Taking $T=r$ and dividing by $N_Q$ proves \eqref{eq:mean-trivial}. For \eqref{eq:mean-pv}, primitive Gauss-sum Fourier inversion and
$\abs{1-e^{2\pi\ii a/r}}\ge4\norm{a/r}_{\R/\Z}$ give the elementary P\'olya--Vinogradov bound
\[
 T\le\sqrt r\,h_{\lfloor r/2\rfloor}
 \le\sqrt r(1+\log r).
\]
The fixed invertible change of residue induced by $4j+3$ has no effect on that bound. Substituting this $T$ completes the proof. No ramified value of a character has been replaced by a unit-modulus value.
\end{proof}

\begin{theorem}[Uniform weighted finite-family comparison]\label{thm:field-average}
For $X\ge1$, $Q\ge2^{16}$, and every common complex array $\abs{u_n}\le1$,
\begin{align}
 &\left|\frac1{N_Q}\sum_{d\in\mathcal D(Q)}
       \sum_{n\le X}b_{-d}(n)u_n
       -\sum_{n\le X}\mu_6(n)u_n\right|
 \le\frac{48X^{3/2}}{\sqrt Q},\label{eq:field-average-weak}\\
 &\left|\frac1{N_Q}\sum_{d\in\mathcal D(Q)}
       \sum_{n\le X}b_{-d}(n)u_n
       -\sum_{n\le X}\mu_6(n)u_n\right|
 \le\frac{96X^{5/4}\sqrt{1+\log X}}{\sqrt Q}.
 \label{eq:field-average}
\end{align}
\end{theorem}
\begin{proof}
For squarefree $n$,
\[
 \prod_{p\mid n}(1+\chi_{-d}(p))=\sum_{r\mid n}\chi_{-d}(r).
\]
The term $r=1$ is exactly the common M\"obius term in the theorem. It is canceled algebraically before taking absolute values. Every other divisor is squarefree and odd, so \eqref{eq:mean-trivial} bounds the remaining expression by
\[
 \frac{24}{\sqrt Q}\sum_{n\le X}\sum_{\substack{r\mid n\\r>1}}\sqrt r
 \le\frac{24X}{\sqrt Q}\sum_{r\le X}r^{-1/2}
 \le\frac{48X^{3/2}}{\sqrt Q}.
\]
Using \eqref{eq:mean-pv} instead gives
\[
 \frac{24X\sqrt{1+\log X}}{\sqrt Q}\sum_{r\le X}r^{-3/4}
 \le\frac{96X^{5/4}\sqrt{1+\log X}}{\sqrt Q}.
\]
Squarefreeness and coprimality were relaxed only in these positive upper sums, not in the original observable.
\end{proof}

Theorem~\ref{thm:field-average} is a quantitative arithmetic comparison, not a cancellation estimate for its first term. If $\theta>1/2$ and $Q=X^q$ with $q\ge5/2-2\theta$, its error fits $X^{\theta+\eps}$. Thus a signed average bound at that power would imply the same bound for $\sum_{n\le X}\mu_6(n)$, and hence ordinary zero exclusion to the right of $\theta$ by Abel summation. The schedules $Q=X^{5/6}$ and $Q=X^{3/2}$ correspond, respectively, to the still-open $5/6$ and square-root targets. They are not achieved zero-free boundaries.

The order of this deduction matters. The finite-family comparison is first applied at each $X$ to recover a \emph{fixed} coefficient sequence. One does not take a Mellin transform while treating a family whose size depends on $X$ as though it were a fixed Dirichlet series. Similarly, arbitrary signed combinations of individual reciprocal residues need not preserve those residues. The theorem supplies a specific coefficient comparison, not a blanket noncancellation assertion.

Finally, deletion of an exceptional set must be paid. Uniformly in $d$,
\[
 \left|\sum_{n\le X}b_{-d}(n)u_n\right|
 \le\sum_{n\le X}2^{\omega(n)}\le X(1+\log X).
\]
Removing $E$ rows while retaining normalization $N_Q$ therefore costs at most $(E/N_Q)X(1+\log X)$. At $Q=X^{5/6}$, an exceptional proportion $Q^{-\delta}$ requires $\delta\ge1/5$ merely to fit an $X^{5/6+\eps}$ allowance. An assertion about almost every auxiliary factor is not a bound for the shared ordinary source.
